% ---- pg-247 （书页 235） ----

益的．在本例中，$\sum\limits_{l=1}^{3}n_l=n=516$ 很大，所以 $H_0$ 可在 $\alpha=0.01$ 显著性水平下进行检验，只要比较
\[
-\left[n-1-(p+g)/2\right]\ln\left(\frac{|\mathbf{W}|}{|\mathbf{B}+\mathbf{W}|}\right)=-511.51\ln(0.771\,4)=132.76
\]
与 $\chi_{p(g-1)}^{2}(0.01)=\chi_{8}^{2}(0.01)=20.09$ 的大小即可．由于 $132.76>\chi_{8}^{2}(0.01)=20.09$，故拒绝 $H_0$．这与刚才用 $F$ 统计量所得出的结果相一致．\hfill$\blacksquare$

\section{处理效应的联合置信区间}
\phantomsection\label{sec:6.5}

当处理效应的相等性假设被拒绝时，人们自然就会对导致假设被拒绝的那些效应感兴趣．对成对比较，用庞弗罗尼方法（见 5.4 节）可构造出 $\boldsymbol{\tau}_k-\boldsymbol{\tau}_l$（或 $\boldsymbol{\mu}_k-\boldsymbol{\mu}_l$）的分量的联合置信区间，这些区间比用别的方法构造的区间要窄，而且只要求用一元 $t$ 分布的临界值．

设 $\tau_{ki}$ 为 $\boldsymbol{\tau}_k$ 的第 $i$ 个分量，由于 $\hat{\boldsymbol{\tau}}_k=\bar{\mathbf{x}}_k-\bar{\mathbf{x}}$ 为 $\boldsymbol{\tau}_k$ 的估计量，故有
\phantomsection\label{eq:6-45}
\begin{equation}
\hat{\tau}_{ki}=\bar{x}_{ki}-\bar{x}_i \tag{6-45}
\end{equation}

而且 $\hat{\tau}_{ki}-\hat{\tau}_{li}=\bar{x}_{ki}-\bar{x}_{li}$ 为两独立样本的均值差．对经过适当修正的 $\alpha$ 值，建立在二样本 $t$ 分布基础上的置信区间依然成立．注意：
\[
\mathrm{Var}(\hat{\tau}_{ki}-\hat{\tau}_{li})=\mathrm{Var}(\bar{X}_{ki}-\bar{X}_{li})=\left(\frac{1}{n_k}+\frac{1}{n_l}\right)\sigma_{ii}
\]

其中 $\sigma_{ii}$ 为 $\boldsymbol{\Sigma}$ 的第 $i$ 个对角元素，正如\hyperref[eq:6-41]{式(6-41)}所暗示的，$\mathrm{Var}(\bar{X}_{ki}-\bar{X}_{li})$ 可通过将 $\mathbf{W}$ 的相应元素除以其自由度来估计，即
\[
\widehat{\mathrm{Var}}\;(\bar{X}_{ki}-\bar{X}_{li})=\left(\frac{1}{n_k}+\frac{1}{n_l}\right)\frac{w_{ii}}{n-g}
\]

此处 $w_{ii}$ 为 $\mathbf{W}$ 的第 $i$ 个对角元素，且 $n=n_1+\cdots+n_g$．

尚需将误差率分配到许多置信区间中去．我们仍可应用\hyperref[eq:5-28]{式(5-28)}．由于有 $p$ 个变量及 $g(g-1)/2$ 对差值，故每一个二样本 $t$ 区间需用到临界值 $t_{n-g}(\alpha/2m)$，其中
\phantomsection\label{eq:6-46}
\begin{equation}
m=pg(g-1)/2 \tag{6-46}
\end{equation}
为联合置信区间的个数．

\mvthm{结论 6.5}{设 $n=\sum\limits_{k=1}^{g}n_k$．对于模型(6-38)，对所有的分量 $i=1,\cdots p$ 及所有差值 $l<k=1,\cdots,g$，以至少为 $1-\alpha$ 的置信度有}
\[
\tau_{ki}-\tau_{li}\ \text{属于}\ \bar{x}_{ki}-\bar{x}_{li}\pm t_{n-g}\left(\frac{\alpha}{pg(g-1)}\right)\sqrt{\frac{w_{ii}}{n-g}\left(\frac{1}{n_k}+\frac{1}{n_l}\right)}
\]
其中的 $w_{ii}$ 为 $\mathbf{W}$ 的第 $i$ 个对角元素．

下面将利用例 6.10 中的小型疗养所数据来说明如何构造处理均值间成对差的联合区间估计．

\noindent\textbf{例 6.11}\quad （处理间均值差的联合置信区间——小型疗养所）

在例 6.10 中我们看到，小型疗养所的平均费用因所有制类型的不同而异．我们可利用结论 6.5 来估计这些差异的大小．为比较私营疗养所与政府疗养所的设备操作和维护人工费用 $X_3$ 之间的差异，我们可以估计 $\tau_{13}-\tau_{33}$．利用\hyperref[eq:6-39]{式(6-39)}及例 6.10 中的信息，有
